\documentclass{article}
\usepackage{amsthm}
\newtheorem{lemma}{Lemma}
\title{A finite sum convention}
\begin{document}
\begin{abstract}
We record the triangular sum identity, including its empty case.
\end{abstract}
\section{Notation}
For $n\in\{0,1,2,\ldots\}$ define $\tau(n)=n(n+1)/2$.
An empty sum is zero.
\section{Identity}
\begin{lemma}\label{lem:sum}
For each $n\geq0$ one has $\sum_{j=1}^n j=\tau(n)$.
\end{lemma}
\begin{proof}
Both sides are zero at $n=0$. Assuming the formula at $n$,
adding $n+1$ gives $(n+1)(n+2)/2=\tau(n+1)$.
\end{proof}
\end{document}
